\begin{center}
\textsc{Gauss} an \textsc{Olbers}.\quad Göttingen, 1.~März 1827.
\end{center}

\par\noindent\hspace*{\parindent}\mbox{.\ .\ .\ .\ .\ .}\quad Meine Abhandlung oder, genauer gesagt, vielleicht meine erste Abhandlung über die gekrümmten Flächen habe ich abgeschlossen. Ich werde sie der Sozietät aber noch nicht übergeben, da zur Ostermesse ohnehin kein Band
\par\smallskip\noindent\hspace*{0.03\linewidth}{\footnotesize\textbf{IX.}}\hfill{\footnotesize48}\hspace*{0.14\linewidth}\par
\clearpage
\noindent\makebox[0pt][l]{\fontsize{13}{16}\selectfont 378}\hfill{\small\textbf{BRIEFWECHSEL MIT OLBERS.}}\hfill\null\par\medskip
\noindent herauskommt.
Mit dem Druck der beiden von mir 1825 und 1826 eingereichten Abhandlungen über die biquadratischen Reste und Suppl.\ theor.\ combin.\ observ. wurde noch nicht begonnen.
Jene Abhandlung enthält für die \textls[120]{unmittelbare} Verwendung in meinem künftigen Werk über die Messung eigentlich nur ein paar Sätze, nämlich 1)~was zur Berechnung des Exzesses der Summe der 3 Winkel über \(180^{0}\) in einem Dreieck auf einer \textls[120]{nicht} sphärischen Fläche, dessen Seiten kürzeste Linien sind, erforderlich ist, und 2)~wie in diesem Fall der Exzess \textls[120]{ungleich} auf die drei Winkel verteilt werden muss, damit die Sinus proportional zu den gegenüberliegenden Seiten werden.
In der Praxis ist dies zwar ganz unwichtig, weil diese Ungleichheit in der Verteilung selbst bei den größten Dreiecken, die sich auf der Erde messen lassen, tatsächlich nicht wahrnehmbar ist. Die Würde der Wissenschaft verlangt aber, dass man die Natur dieser Ungleichheit klar begreift.
Und so kann man hier, wie so oft, durchaus ausrufen »Tantae molis erat!«, um dahin zu gelangen. ---
Wichtiger als die Lösung dieser 2 Aufgaben ist aber, dass die Abhandlung mehrere allgemeine Prinzipien begründet, aus denen später in einer spezielleren Untersuchung die Lösung vieler wichtiger Aufgaben abgeleitet werden kann.

\par\noindent\hspace*{\parindent}Ich komme noch einmal auf den zu Beginn dieses Briefes erwähnten Gegenstand zurück.
Ich habe allen meinen Hilfstafeln und Rechnungen \textsc{Walbecks} Abplattung von \(\frac{1}{302{,}78}\) zugrunde gelegt. Ich glaube aber, dass alle bisherigen Gradmessungen, wenn man ihre Daten \textls[120]{vollständig} berücksichtigte, zeigen würden, dass \textsc{Sabine}s Abplattung von \(\scriptstyle\frac{1}{2\mkern5mu 8\mkern5mu 9}\) fast ebenso gut damit vereinbar wäre, und mir scheint, dass alle bisher vorliegenden Gradmessungen noch viel zu wenig ausgedehnt sind, um die Abplattung zwischen engeren Grenzen einzuschließen.
Gewiss ist, dass die Erde ein unregelmäßiger Körper ist. Die Polhöhen der Orte (ganz abgesehen von Beobachtungsfehlern) schwanken immer mehrere Bogensekunden (vielleicht hier und da recht viele Bogensekunden) um die Werte, die man unter der Annahme irgendeines regelmäßigen Ellipsoids \textls[120]{berechnet}, und ebenso schwanken die wirklichen Pendellängen um die berechneten (die \textls[120]{mittlere} Schwankung der Pendellängen beträgt vielleicht \(\scriptstyle\frac{1}{3\mkern7mu 0}\) englische Linie).
Aber eben deshalb können Gradmessungen von geringer Ausdehnung wenig zur Kenntnis der mittleren Gestalt der Erde beitragen. Insbesondere halte ich den lappländischen
\par\clearpage
\noindent\null\hfill{\small\textbf{ZUR HANNOVERSCHEN TRIANGULATION.}}\hfill\makebox[0pt][r]{\fontsize{13}{16}\selectfont \textbf{379}}\par\medskip
\noindent Bogen für viel zu klein.
Umso wichtiger scheint mir, dass nach und nach alle präzisen Triangulationen in Europa in einen einzigen Zusammenhang gebracht werden.
Ich habe jetzt die Hoffnung, dass mir die bayerischen Dreiecke mitgeteilt werden.
Wenn erst alle europäischen Sternwarten von Åbo bis Palermo und von Nicolajef bis Dublin durch Dreiecke verbunden sind, sodass ihre Lage zueinander mit aller Genauigkeit bestimmt ist, die die feinsten geodätischen Messungen ermöglichen, dann wird man, d.~i.\ dann werden unsere Nachkommen, alles mit viel größerer Sicherheit beurteilen \mbox{können.\ .\ .\ .\ .\ .\ .}
